\section{Introduction}

A model's weight tensors predict its generalization gap --- how much training accuracy exceeds test accuracy
at convergence --- better than they predict its test accuracy outright. In the Unterthiner CIFAR-10 CNN zoo,
a simple position-indexed MLP achieves Spearman rank correlation $r = 0.56$ on generalization gap prediction
but only $r = 0.28$ on test accuracy prediction, despite receiving the identical weight tensor as input.
The harder prediction target, as it turns out, is the more legible one.

This asymmetry is not merely a curiosity. Generalization gap (train\_acc $-$ test\_acc at convergence) is the
direct quantitative signature of overfitting --- the quantity that determines whether a model has memorized
its training data or extracted transferable patterns. For practitioners managing large model zoos, identifying
low-gap models without held-out test evaluation would dramatically reduce the cost of model selection. For
theorists, knowing that weight tensors encode an accessible gap signal raises a structural question: what
is the geometric nature of overfitting in weight space, and which architectural inductive biases are best
suited to extract it?

\textbf{The gap in existing work.} A rich literature has established that weight-space encoders can predict test
accuracy with Spearman correlations as high as $r \approx 0.9$. \citet{unterthiner2020} introduced the model zoo
benchmark and flat MLP baseline. \citet{navon2023} demonstrated that permutation-equivariant Deep Weight
Space networks (DWS) substantially outperform flat baselines on test accuracy. \citet{zhou2023} showed
that Neural Functional Transformers (NFT) --- cross-layer attention architectures respecting weight-space
symmetries --- achieve competitive or superior performance. \citet{kofinas2024} extended this with graph
neural network encoders over the neural network graph. Critically, \textbf{all these works benchmark exclusively
on test accuracy prediction.} Generalization gap is computable from the same zoo metadata yet has never
been treated as the primary prediction target in this line of work --- and no study has examined whether
encoder architecture advantage is target-specific (gap vs.\ test\_acc) or uniform across targets.

\textbf{Our approach.} We conduct the first controlled comparison of equivariant vs.\ non-equivariant weight
encoders on generalization gap as the primary prediction target. The central methodological innovation is
a dual-target experimental design: the same encoder trained on the same zoo with the same hyperparameter
budget, evaluated on both gap and test accuracy prediction. This design isolates target-specificity from
encoder-specificity. We additionally introduce a partial Spearman analysis --- the first in this literature ---
that tests whether gap predictions contain information about true gap that is statistically independent of
test accuracy rank.

\textbf{Key finding.} NFT's cross-layer attention captures gap-specific overfitting structure in weight tensors
that is statistically independent of test accuracy: partial Spearman(NFT\_gap\_pred, true\_gap $|$ true\_test\_acc)
$= 0.73$, $p = 1.6 \times 10^{-167}$. This result --- substantially exceeding NFT's direct gap Spearman of 0.57 ---
indicates that the gap-specific component of NFT's predictions is stronger than the combined signal, and
that generalization gap and test accuracy occupy distinct regions of the weight-space information landscape.

\textbf{Contributions.} We make four contributions:

\begin{enumerate}
\item \textbf{Gap learnability (existence).} We establish that generalization gap is predictable from weight tensors
   at Spearman $r > 0.5$ using both position-indexed (FlatMLP $r = 0.5567$) and equivariant (DWSNet $r = 0.5104$)
   encoders, with a confirmed A1 audit showing gap is not trivially derivable from test accuracy
   (Spearman(gap, $-$test\_acc) $= -0.142$).

\item \textbf{Architecture specificity.} Among the four tested encoders, NFT achieves the highest gap Spearman
   ($r = 0.5752$, 95\% CI: [0.5339, 0.6158]). FlatMLP's 95\% CI from the same run is [0.4850, 0.5801];
   the two intervals overlap at the margin ([0.5339, 0.5801]), but NFT's lower bound exceeds FlatMLP's
   point estimate (0.5330), providing directional statistical evidence of advantage. DWSNet ($r = 0.4881$)
   and GNN ($r = 0.3747$) both underperform FlatMLP on gap --- showing that equivariance alone is
   insufficient; cross-layer reasoning matters.

\item \textbf{Gap-specific signal independence (P3).} NFT gap predictions contain substantial information about
   true generalization gap that is independent of test accuracy (partial Spearman $r = 0.73$, $p \approx 0$),
   establishing gap as a genuinely distinct weight-space prediction target.

\item \textbf{Documented null result ($\Delta$ mechanism).} The differential advantage hypothesis --- that equivariant
   encoders show a larger Spearman improvement on gap than on test\_acc ($\Delta > 0.02$) --- is not confirmed
   under our experimental conditions. We identify a plausible confounder (anomalously low FlatMLP test\_acc
   in our zoo, $r = 0.28$ vs.\ literature $\approx 0.85$) and propose corrected future experiments.
\end{enumerate}

We organize the paper as follows. Section~\ref{sec:related} surveys weight-space encoder and generalization prediction
literature, positioning our dual-target formulation against prior work. Section~\ref{sec:method} describes the experimental
methodology including the dual-target design and partial Spearman analysis. Section~\ref{sec:setup} presents experimental
setup. Section~\ref{sec:results} presents results. Section~\ref{sec:discussion} discusses mechanism interpretation, limitations, and broader
implications. Section~\ref{sec:conclusion} concludes.
